\section{Conclusion}\label{sec:conclusion}

This survey began from a simple distinction: an agent that carries state across tasks is not the same kind of system as one that resets. The episodic assistant that answers a question and forgets it can be wrong without making the error durable. The always-on agent inherits what it wrote yesterday. Its competence and its liabilities both compound, and the substrate that lets it carry knowledge forward is the same substrate that lets a stale preference, a poisoned trace, or a lapsed permission govern an action taken months later. Such a system should be analyzed as \emph{persistent state}, not memory alone: retrievable memory plus task ledgers, permissions, tool and credential state, standing commitments, provenance records, shared and social state, and trigger conditions, with external commitments treated as part of the accountability surface because they must remain traceable to the state that authorized them. Memory is the most studied component of that union, but treating the whole as memory hides the governance obligations that make always-on agents distinct.

\subsection{What the survey established}

Four contributions structure the argument. First, always-on agents are persistent-state systems. Their state is governed over a lifecycle, observe, write, validate, organize, retrieve, act, update, forget, audit, and rollback, and can be characterized along six axes: authority (who may write or invoke it), scope (where it may legitimately apply), mutability (how it changes under new evidence), provenance (where it came from and through what transformations), recoverability (whether a change can be undone), and actionability (whether it can drive an external effect). Five invariants tie the axes to the lifecycle: authority monotonicity, scope non-expansion, deletion propagation, provenance preservation, and rollback traceability. The frame is wider than memory by construction. Classical cognitive architectures already separated procedural, episodic, and semantic stores within a single decision cycle \citep{laird2012soar, anderson2004actr, tulving1972episodic}, and the complementary-learning-systems tradition explained why fast capture and slow consolidation must be kept on different timescales to avoid catastrophic interference \citep{mcclelland1995complementary}. Modern agent frameworks ported these modules into LLM systems \citep{sumers2024cognitivearchitectureslanguageagents}, but they ported the storage without the governance: the cognitive theories assumed type boundaries and decay schedules enforced by biology, whereas an LLM agent mixes fast and slow updates in flat text with no runtime enforcement of who wrote what, with what authority, recoverable by whom.

Second, the literature is read as a layered persistent-state stack rather than as a memory taxonomy, so that substrates, the lifecycle, continual adaptation, governance, multi-agent state, evaluation, failure modes, and application domains each receive a first-class treatment. This is not a cosmetic relabeling. A memory-form taxonomy strands governance, serving substrates, deletion propagation, and rollback as footnotes, which is precisely the failure mode the survey exists to name. The stack instead lets the memory-form and lifecycle axes live where they are sharpest, inside the substrate and state-movement layers, while giving authority, provenance, and recoverability the room to carry the argument.

Third, we coded a $435$-work corpus and measured which parts of the lifecycle each work addresses. The main finding is a structural asymmetry between accumulation and governance. In our coding, the forward arc of the lifecycle is crowded: retrieve appears in $269$ works and write in $200$. The return arc is sparse: only $27$ of $435$ works expose any rollback mechanism, and across the state axes authority is the rarest at $72$ of $435$. These counts are scoping estimates, but they consistently show that the works we examined study storing and retrieving more thoroughly than revoking, recovering, attributing, and bounding stored state.

Fourth, we turned the diagnosis into a checkable protocol by proposing the Always-On Evaluation Protocol and instantiating it as AOEP-v0: an event-stream and snapshot schema, worked episodes, a validator, and a pilot that scores state mutation and recovery rather than answer quality. The pilot is deliberately small and representational rather than a benchmark. It illustrates how tested memory wrappers can fail obligations that require explicit governance state, such as permission epochs, deletion ledgers, trust tiers, conflicts, and rollback records.

\subsection{Why bigger context and better retrieval are not enough}

One objection to our thesis is that the problem is already dissolving on its own, that million-token context windows and stronger retrieval will let an agent read all of its past and behave correctly. The corpus does not support this. Larger context and retrieval address a reading problem, and the governance gaps we document are not reading problems. Long-context studies show that capacity is not comprehension: models underuse information placed in the middle of a long window even when it is present \citep{liu2023lostmiddlelanguagemodels}, and retrieval quality decomposes into noise robustness, rejection, and faithfulness failures that more tokens do not fix \citep{chen2023benchmarkinglargelanguagemodels}. Even granting perfect recall, retrieval answers ``what did I record?'' It does not answer the four questions an always-on agent must answer to act safely: what \emph{should} persist, which competing record is \emph{authoritative}, how a record should \emph{change} when new evidence contradicts it, and how a user can \emph{revoke} it. These are decisions about authority, mutability, and recoverability, not about context length.

The difference becomes sharp where state crosses from passive storage into action. A preference that was true last month and was changed yesterday is still retrievable today; the question is whether the agent treats it as superseded or as current, and that is a mutability-and-provenance decision no retriever makes \citep{uddin2026recallforgettingbenchmarkinglongterm, ma2026memprobe}. Personalization research has repeatedly shown that accumulated user state drifts, contradicts itself, and can be poisoned: routine interactions gradually weaken confirmation boundaries and expand the agent's effective action scope without any single step looking like an attack \citep{xu2026toxicchats}. Surveys of personalized agents catalog these requirements, multi-session coherence, preference retrieval, drift handling, while noting that the systems offer no authority monotonicity and no scope gating to enforce them \citep{xu2026personalizedllmpoweredagentsfoundations}. A bigger window lets the agent see the poisoned preference more reliably. It does not give the agent or the user any mechanism to mark it untrusted, bound where it applies, or roll back the actions it already licensed. Recent state-trajectory work makes the same point formally: stores that keep conclusions while dropping sources become hard to correct \citep{kwon2026reclaim}, and record-level stores cannot satisfy rollback traceability without additional governed-state structure \citep{orogat2026gem}. The deficiency is representational and operational; scale alone does not remove it.

\subsection{The operational agenda}

The corrective is to treat persistent state as a measured, governed lifecycle rather than as a passive log to be searched. Concretely, this means three commitments that recur across every part of this survey and that we summarize in Table~\ref{tab:conclusion-shift}. State decisions must be made explicitly rather than implicitly: a system should decide what to admit at the write boundary, what authority a record carries, and under what scope it may act, instead of accumulating everything and hoping retrieval sorts it out. State must remain attributable and recoverable: provenance must survive consolidation so that a corrupted or superseded record can be located, quarantined, and undone, rather than dropped the moment a summary scores well on aggregate recall. And state governance must be evaluated rather than asserted: a system that claims to honor deletion, respect authority, or support rollback should be driven through perturbations that exercise those obligations and scored on whether the obligations actually hold.

\begin{table}[t]
\centering
\caption{The shift the survey argues for: from persistent state as an accumulated log to persistent state as a governed lifecycle. Each row pairs the prevailing framing with the operational alternative and the corpus signal that motivates it.}
\label{tab:conclusion-shift}
\small
\begin{tabular}{@{}p{0.28\columnwidth}p{0.34\columnwidth}p{0.28\columnwidth}@{}}
\toprule
\textbf{Question} & \textbf{Accumulate-and-retrieve framing} & \textbf{Governed-lifecycle alternative} \\
\midrule
What persists? & Store the interaction, retrieve later. & Admit at the write boundary under a scope and authority label. \\
Which record rules? & Most similar to the query. & The authoritative one; authority only narrows. \\
How does it change? & Append a newer entry. & Mutate with provenance; mark prior truth superseded. \\
How is it revoked? & Delete the row, hope it propagates. & Cascade deletion, verify completeness, log the rollback. \\
How do we know? & Downstream task accuracy. & Score the governance obligation directly. \\
\bottomrule
\end{tabular}
\end{table}

These commitments are practical because isolated pieces of each already exist in the corpus. Operating-system framings externalize working state to durable storage and reason about what should be paged in and out, supplying the substrate discipline a lifecycle needs \citep{packer2024memgptllmsoperatingsystems}. Provenance-rooted and immutable-ledger designs show that source attribution and tamper-evidence can be maintained through updates \citep{wright2025immutablemem}, and certified defenses demonstrate that write-time provenance plus randomized ablation can bound the influence of poisoned writes \citep{sharma2026smsrcertified}. Transactional models import compensable execution and rollback from databases into multi-agent planning \citep{chang2025sagallm}. Workflow-abstraction methods show that accumulation can be made selective and reusable rather than indiscriminate \citep{wang2024agentworkflowmemory}. The missing step is integration across all six axes and five invariants, with evaluation as obligations rather than as incidental byproducts of task success. AOEP-v0 starts from that integration by separating positive obligations a system must actively satisfy from negative no-leakage invariants that a system can satisfy trivially by storing nothing. The pilot suggests that obligations reducible to recall and obligations requiring maintained governance state should be reported separately.

\subsection{Outlook}

The stakes rise rather than fall as deployments mature. The failure modes we cataloged compound under continual operation: poisoning persists across sessions once it enters durable state \citep{xie2026crosssessioninjection}, alignment can degrade as the agent accumulates and reuses its own outputs \citep{shao2025misevolution}, and in shared settings evaluator bias in stored trajectories propagates to future agents with no safe contamination threshold even under oracle consolidation \citep{liu2026memorycontagion}. None of these is a single-turn prompt problem, and none is solved by reading more of the past more reliably. They are governance problems in the sense this survey defines: failures of authority, scope, provenance, and recoverability over a lifecycle that the works we coded instrument mostly on its forward arc. We read our main numbers, $27$ of $435$ works with any rollback and $72$ of $435$ touching authority, not as a verdict on the quality of existing work, which is often excellent within its frame, but as a map of where the frame stops. Persistence is not a minor feature of always-on agents to be handled by a larger window or a sharper index. It is the property that makes these systems both more capable and more accountable than their episodic predecessors, and treating it as a measured, governed lifecycle, deciding what should persist, which state is authoritative, how it changes, and how it is revoked, is the operational agenda the next generation of always-on agents will be built and judged against.
